\chapter{Testing and Acceptance}
\label{ch:testing}

This section specifies the complete testing strategy for gert: the test framework, unit and contract test patterns, integration test strategy, the \texttt{gert test} feature, regression testing against the runbook corpus, performance acceptance criteria, the extension test harness, and CI/CD gate requirements.

% ------------------------------------------------------------------ %
\section{Test Framework}
% ------------------------------------------------------------------ %

\subsection{Standard Library and testify}

All gert tests are written in Go using the standard \texttt{testing} package as the test runner and \texttt{github.com/stretchr/testify}~\cite{testify-library} for assertions:

\begin{itemize}
  \item \texttt{testify/assert} — non-fatal assertions; test continues on failure.
  \item \texttt{testify/require} — fatal assertions; test halts on failure.
  \item \texttt{testify/mock} — interface mocking for unit tests (engine, event bus, trace writer).
\end{itemize}

Table-driven tests are the preferred pattern for all logic that has well-defined input/output pairs (e.g.\ expression evaluation, schema validation, governance rule matching):

\begin{minted}{go}
func TestGovernanceAllowlist(t *testing.T) {
    cases := []struct {
        name    string
        cmd     []string
        allowed bool
    }{
        {"exact match",    []string{"kubectl", "get", "pods"}, true},
        {"partial match",  []string{"kubectl", "delete", "ns"}, false},
        {"empty allowlist", []string{"echo", "hi"}, false},
    }
    for _, tc := range cases {
        t.Run(tc.name, func(t *testing.T) {
            result := allowlist.Check(tc.cmd)
            assert.Equal(t, tc.allowed, result)
        })
    }
}
\end{minted}

\subsection{Test Packages}

Each Go package in the module ships a co-located \texttt{\_test.go} file (white-box tests) and optionally an \texttt{<pkg>\_integration\_test.go} file (integration tests, gated by a build tag). The naming convention:

\begin{verbatim}
pkg/engine/engine.go
pkg/engine/engine_test.go           // white-box unit tests
pkg/engine/engine_integration_test.go  // build tag: //go:build integration
\end{verbatim}

% ------------------------------------------------------------------ %
\section{The \texttt{gert test} Feature}
% ------------------------------------------------------------------ %

\subsection{Scenario Replay Testing}

\texttt{gert test <runbook.yaml>} discovers and executes all scenario tests for a runbook. A scenario test consists of:

\begin{itemize}
  \item \textbf{\texttt{inputs.yaml}} — resolved input values for the run.
  \item \textbf{\texttt{scenario.yaml}} (optional) — pre-recorded command responses and manual evidence.
  \item \textbf{\texttt{test.yaml}} — assertions over the resulting trace and captures.
\end{itemize}

Scenario directories are located by convention:

\begin{verbatim}
{runbook-dir}/scenarios/{runbook-name}/{scenario-name}/
  inputs.yaml
  scenario.yaml   # optional; omit for runbooks with no CLI/tool steps
  test.yaml
\end{verbatim}

The test runner executes the runbook in replay mode using the scenario file as the command/evidence fixture, then evaluates assertions from \texttt{test.yaml}.

\subsection{Test Assertions (\texttt{test.yaml})}

\begin{minted}{yaml}
assertions:
  outcome: passed           # "passed" | "failed" | "skipped"
  steps_passed: 5
  steps_failed: 0

  captures:
    db_host:
      equals: "prod-db.example.com"
    migration_status:
      contains: "complete"

  trace:
    - type: step/completed
      step_id: run-migration
    - type: governance/blocked
      step_id: drop-table
      data.rule: denylist
\end{minted}

The \texttt{trace} assertions match events in order of occurrence. The \texttt{data.*} dotted paths dereference into the event's \texttt{data} object.

% ------------------------------------------------------------------ %
\section{Contract Tests}
% ------------------------------------------------------------------ %

Contract tests verify the interface boundaries between the schema's five components. Each contract test lives in the \emph{consumer} package and uses a conformance helper that validates any implementation satisfying the interface.

\subsection{Core Runtime Contract}

The runtime \texttt{Executor} interface must satisfy:

\begin{itemize}
  \item Emits a \texttt{run/started} event before executing any step.
  \item Emits a \texttt{step/started} before and \texttt{step/completed}/\texttt{step/failed} after each step.
  \item Emits \texttt{run/completed} or \texttt{run/failed} as the final event.
  \item Does not emit events after the terminal event.
  \item Does not call any adapter method directly (events only).
\end{itemize}

\begin{minted}{go}
// ConformanceTest validates any Runtime implementation.
func ConformanceTest(t *testing.T, rt Runtime) {
    t.Helper()
    bus := newTestEventBus()
    err := rt.Execute(context.Background(), fixtureRunbook, bus)
    require.NoError(t, err)
    events := bus.Events()
    require.Equal(t, "run/started",   events[0].Type)
    require.Equal(t, "run/completed", events[len(events)-1].Type)
    assertNoEventsAfterTerminal(t, events)
}
\end{minted}

\subsection{Schema Contract Tests}

\begin{itemize}
  \item Every \texttt{apiVersion: runbook/v1} document in \texttt{testdata/} must pass structural validation without errors.
  \item Every invalid fixture in \texttt{testdata/invalid/} must fail validation with a structured error (message + path).
  \item Structural and semantic validation are tested independently: a structurally valid document may fail semantic validation (e.g.\ referencing an undefined capture).
\end{itemize}

\subsection{Extension Manifest and Handshake Tests}

The extension host verifies:

\begin{itemize}
  \item An extension process that responds to the initialisation handshake within the configured timeout is accepted.
  \item An extension process that fails to respond is terminated and its capabilities are not registered.
  \item Capability grants in the manifest are enforced: an extension without \texttt{cap:schema.register} cannot register schema.
  \item A manifest with missing required fields is rejected at load time with a structured error.
\end{itemize}

\subsection{Tool Invocation Envelope Tests}

\begin{itemize}
  \item Tool invocations carry the correct \texttt{run\_id} and \texttt{step\_id} correlation fields.
  \item Tool outputs are captured and available as step captures after the step completes.
  \item A tool that exits with non-zero triggers \texttt{step/failed} with \texttt{error\_type: tool\_error}.
  \item Timeout cancellation: a tool that exceeds its timeout receives a \texttt{context.DeadlineExceeded} and emits \texttt{step/failed} with \texttt{error\_type: timeout}.
\end{itemize}

\subsection{Event Bus Determinism Tests}

\begin{itemize}
  \item Events delivered to the event bus maintain the \texttt{seq} ordering guaranteed by \texttt{RunStore.WriteTrace}.
  \item Slow adapter consumers do not cause the engine's execution loop to block (bounded channel or dropping policy required).
  \item Event replay (\texttt{ReadTraceSince}) returns events in ascending \texttt{seq} order.
\end{itemize}

% ------------------------------------------------------------------ %
\section{Integration Test Strategy}
% ------------------------------------------------------------------ %

\subsection{End-to-End Fixture Runs}

Integration tests spin up the gert runtime against fixture runbooks stored in \texttt{testdata/runbooks/} and assert on the resulting trace files. They are gated by the \texttt{integration} build tag and run in CI separately from unit tests.

\begin{minted}{go}
//go:build integration

func TestE2E_BranchingRunbook(t *testing.T) {
    dir := t.TempDir()
    runID, err := exec.RunFixture(dir, "testdata/runbooks/branch.runbook.yaml",
        WithInputs(map[string]string{"env": "prod"}),
        WithScenario("testdata/scenarios/branch-prod.yaml"),
    )
    require.NoError(t, err)
    trace := readTrace(t, dir, runID)
    assertEvent(t, trace, "step/completed", map[string]any{
        "data.step_id": "deploy-prod",
    })
    assertNoEvent(t, trace, "step/completed", map[string]any{
        "data.step_id": "deploy-staging",
    })
}
\end{minted}

\subsection{Fixture Runbook Corpus}

The integration test corpus includes fixtures covering:

\begin{itemize}
  \item Linear runbook (all step types: cli, manual, tool, invoke).
  \item Branching runbook (both branch arms; governance-blocked branch).
  \item Iterate runbook (list mode, convergence mode, parallel concurrency).
  \item Nested invoke (parent + two child runbooks).
  \item Retry policy (transient failure followed by success).
  \item Resumption (crash after step 2, resume, verify trace continuity).
  \item Governance: allowlist violation, denylist violation, approval gate.
\end{itemize}

\subsection{Trace Assertion Helpers}

A shared test package \texttt{pkg/testutil} provides helpers:

\begin{minted}{go}
package testutil

// ReadTrace reads and parses all events from a run's trace.jsonl.
func ReadTrace(t *testing.T, runDir, runID string) []engine.DurableEvent

// AssertEvent asserts that at least one event matches type and data predicates.
func AssertEvent(t *testing.T, events []engine.DurableEvent,
    eventType string, predicates map[string]any)

// AssertEventOrder asserts events appear in the given type order.
func AssertEventOrder(t *testing.T, events []engine.DurableEvent,
    types ...string)

// AssertNoEvent asserts no event matches the given type and predicates.
func AssertNoEvent(t *testing.T, events []engine.DurableEvent,
    eventType string, predicates map[string]any)
\end{minted}

% ------------------------------------------------------------------ %
\section{Regression Testing}
% ------------------------------------------------------------------ %

\subsection{Runbook Corpus as Fixtures}

The runbook corpus (located in \texttt{examples/} and operator-contributed runbooks) serves as a regression suite for schema stability. The regression test runner:

\begin{enumerate}
  \item Reads each runbook from the corpus.
  \item Runs \texttt{gert validate} (schema) and asserts it passes.
  \item For runbooks with associated scenario files, runs \texttt{gert test} and asserts the outcome matches the golden output.
  \item Fails the build if any previously passing corpus runbook regresses.
\end{enumerate}

The schema acceptance target is $\geq$95\% of corpus runbooks passing \texttt{gert validate} unmodified between minor releases.

\subsection{Golden Trace Comparisons}

For deterministic scenarios, the expected trace can be stored as a golden file. The golden comparison test:

\begin{enumerate}
  \item Runs the scenario.
  \item Normalises timestamps and run IDs (replaces with fixed values).
  \item Compares the normalised JSON-lines output to \texttt{testdata/golden/<scenario>.jsonl}.
  \item Fails if any event differs. Golden files are updated by running with \texttt{-update} flag.
\end{enumerate}

% ------------------------------------------------------------------ %
\section{Performance Acceptance Criteria}
% ------------------------------------------------------------------ %

The following latency budgets must pass on a reference machine (4-core, 8 GB RAM, SSD) with the integration test suite:

\begin{center}
\begin{tabular}{lll}
\hline
\textbf{Operation} & \textbf{p50 target} & \textbf{p99 target} \\
\hline
\texttt{gert validate} (single runbook, <50 steps) & $<$50 ms   & $<$200 ms  \\
Step startup (cli step)                              & $<$10 ms   & $<$50 ms   \\
Event delivery to adapter (local channel)            & $<$1 ms    & $<$5 ms    \\
Trace event write + fsync                            & $<$5 ms    & $<$20 ms   \\
Snapshot write (100 steps history)                   & $<$20 ms   & $<$80 ms   \\
\texttt{gert test} (10-step replay scenario)         & $<$500 ms  & $<$1 s     \\
Tool invocation round-trip (stdio, local process)    & $<$50 ms   & $<$200 ms  \\
\hline
\end{tabular}
\end{center}

Performance benchmarks are written using Go's \texttt{testing.B} and run in CI weekly (not on every PR). A regression of $>$2$\times$ on any metric blocks the nightly build.

% ------------------------------------------------------------------ %
\section{Extension Test Harness}
% ------------------------------------------------------------------ %

\subsection{Overview}

Extension developers can test their extension logic without spinning up the full gert runtime. The \texttt{pkg/exttest} package provides a test harness:

\begin{minted}{go}
import "github.com/ormasoftchile/gert/pkg/exttest"

func TestMyExtension(t *testing.T) {
    h := exttest.NewHarness(t,
        exttest.WithManifest("testdata/my-ext.manifest.yaml"),
        exttest.WithCapabilities("cap:schema.register", "cap:tool.register"),
    )
    defer h.Close()

    // Start the extension process
    h.Start("./my-extension")

    // Assert it registered expected tools
    tools := h.RegisteredTools()
    require.Contains(t, tools, "my-tool")

    // Invoke a tool and assert the response
    resp := h.InvokeTool("my-tool", map[string]any{"arg1": "value"})
    assert.Equal(t, 0, resp.ExitCode)
    assert.Contains(t, resp.Stdout, "expected output")
}
\end{minted}

\subsection{Harness Capabilities}

The test harness provides:

\begin{itemize}
  \item A mock extension host that implements the full handshake protocol.
  \item Assertion helpers for registered tools, schema registrations, and event subscriptions.
  \item A \texttt{FakeRunContext} that provides a realistic run context without executing a runbook.
  \item Automatic process cleanup on test completion.
  \item Configurable capability grants (to test capability enforcement).
\end{itemize}

\subsection{Mocking the Extension Transport}

For pure unit tests of extension logic (without spawning a subprocess), the transport layer is mockable:

\begin{minted}{go}
// InProcessTransport runs the extension handler in-process for unit testing.
type InProcessTransport struct {
    Handler ExtensionHandler
}
\end{minted}

% ------------------------------------------------------------------ %
\section{Acceptance Criteria}
% ------------------------------------------------------------------ %

The following acceptance criteria must be satisfied before a release is considered shippable:

\begin{enumerate}[label=\arabic*.]
  \item Core runtime has no direct adapter dependencies — verified by the \texttt{no\_import} contract test.
  \item Extensions register schema and tools through declared contracts — verified by the extension handshake conformance suite.
  \item Capability violations fail fast with structured errors — verified by the governance contract test for each capability type.
  \item Adapters integrate using the same stable core contracts — verified by running each adapter's integration test against a mock core.
  \item Every step type (\texttt{cli}, \texttt{manual}, \texttt{tool}, \texttt{invoke}, \texttt{branch}, \texttt{iterate}) has $\geq$1 end-to-end integration test covering the happy path and at least one failure path.
  \item The runbook corpus passes \texttt{gert validate} at $\geq$95\% without modification.
  \item All performance targets in \S~8.6 pass on the reference machine.
  \item \texttt{gert test} successfully replays all scenario fixtures in \texttt{testdata/scenarios/}.
\end{enumerate}

% ------------------------------------------------------------------ %
\section{CI/CD Integration}
% ------------------------------------------------------------------ %

\subsection{PR Gate (every pull request)}

\begin{enumerate}[label=\arabic*.]
  \item \textbf{Unit tests}: \texttt{go test ./...} (no build tags). Target: $<$60 s.
  \item \textbf{Schema validation}: validate all YAML fixtures in \texttt{testdata/} against JSON Schema. Target: $<$10 s.
  \item \textbf{Contract tests}: \texttt{go test -run Contract ./...}. Target: $<$30 s.
  \item \textbf{Linter}: \texttt{golangci-lint run} with the project's \texttt{.golangci.yaml} config.
  \item \textbf{Vet}: \texttt{go vet ./...}.
\end{enumerate}

All five gates must pass for a PR to be mergeable.

\subsection{Integration Gate (merge to main)}

\begin{enumerate}[label=\arabic*.]
  \item \textbf{Integration tests}: \texttt{go test -tags=integration ./...}. Target: $<$5 min.
  \item \textbf{Regression corpus}: \texttt{gert test} against the corpus. Target: $<$2 min.
  \item \textbf{Golden trace comparison}: normalised trace comparison for all deterministic fixtures.
\end{enumerate}

\subsection{Nightly Gate}

\begin{enumerate}[label=\arabic*.]
  \item \textbf{Performance benchmarks}: \texttt{go test -bench=. -benchtime=10s ./...} — compare against stored baseline; fail on $>$2$\times$ regression.
  \item \textbf{Race detector}: \texttt{go test -race ./...}.
  \item \textbf{Extension conformance suite}: run all built-in extensions through the conformance harness.
\end{enumerate}
